\paragraph{Mesh-based measurement.} Marching cubes \citep{lorensen1987} is the standard way to extract
an iso-surface from a voxel volume; we use the topologically consistent variant of \citet{lewiner2003}
as implemented in scikit-image \citep{vanderwalt2014}. \citet{nystrom2002} compared areas of digital
(voxel-face) surfaces and triangulated surfaces on phantoms and real objects and found the
triangulated areas considerably more accurate, with errors decreasing as the voxel size decreased.

\paragraph{Local estimators and multigrid convergence.} A large body of work in digital geometry
estimates surface area from binary volumes by assigning a weight to each local $2\times2\times2$
configuration. \citet{mullikin1993} derived optimal weights for digitized planes, and
\citet{windreich2003} and \citet{lindblad2005} optimised local configuration weights to minimise the
bias for isotropically oriented surfaces. A central theoretical question is multigrid convergence,
i.e.\ whether the estimate converges to the true area as the grid is refined. \citet{kenmochi2000}
showed that local polyhedrization schemes such as marching cubes on binary data are not multigrid
convergent even for spheres, cubes and cylinders, and \citet{svane2013} proved that no local digital
algorithm for surface area is multigrid convergent on the class of convex polytopes.
\citet{ziegel2010} studied design-based surface-area estimation from digitizations with random grid
placement. \citet{klette2004} give a general account. These results explain why the binary estimators
in this work keep a resolution-independent bias, and they motivate using greyscale information.

\paragraph{Learning-based estimation.} Deep networks operating directly on voxel grids
\citep{maturana2015} and on unordered point sets \citep{qi2017} are established tools for 3D shape
analysis. For CT of porous materials, \citet{rabbani2020} trained a convolutional network to predict a
range of morphological, hydraulic, electrical and mechanical properties of porous samples from 3D
images in a fraction of the time of image-based simulation. Random Forests \citep{breiman2001} on
hand-crafted descriptors remain a strong baseline for small tabular problems.

\paragraph{Gap.} The work above either derives fixed estimators for binary data or learns global
properties of complex microstructures, where an exact ground truth is not available. We did not find a
study that measures how much a learned model reduces surface-area bias relative to marching cubes
against an exact reference, separately for resolution, noise level and unseen shape families, and that
compares deep models with a calibrated feature-based baseline. That comparison is the subject of this
work.
